\subsection{Relation to existing nonvanishing and rigidity theorems}

Kriz--Nordentoft already develop horizontal \(p\)-adic \(L\)-functions,
refined maximum-modulus principles, and simultaneous nonvanishing under
specified arithmetic hypotheses~\cite{kn}. Serban's infinitesimal
multiplicative Manin--Mumford theorem addresses rigidity of torsion
intersections for formal schemes over finite extensions of
\(\Q_p\)~\cite{serban}. These results make the surrounding area an
established research program. The result proved here is the particular
finite-grid, total-positive-loss statement for bounded analytic data;
this analysis does not establish that it is absent from the entire
literature.

The matrix corollary suggests a disciplined route to a genuinely new
arithmetic theorem: identify a specific Selmer presentation or a
determinantal ideal in a \(p\)-adic family, normalize its coefficients
integrally, and produce finite nonzero witnesses with a total budget.
The selector then supplies a simultaneous specialization. The witness
step and the passage from a bounded normalized length to the exact
desired arithmetic assertion must be proved separately.
Section~\ref{sec:counterexample} gives another precise boundary on
overly strong specialization rules.

